\reading{IM-15}{Market-Driven Scenarios: An Approach for Plausible Scenario Construction}{DEFER}{0}

\srcnote{Bennett W. Golub (ed.), \textit{BlackRock's Guide to Fixed-Income Risk
Management} (Wiley, 2024), Chapter 5. No GARP source book covers this reading.}

\begin{imkeybox}[title={The four objectives in this reading}]
\begin{enumerate}[label=\textbf{\alph*.}]
\item Describe and evaluate the market-driven scenario (MDS) approach used to develop scenarios for stress testing investment portfolios.
\item Describe the scenario z-score, the volatility z-score, the correlation z-score, and the Mahalanobis distance, and explain how these are used in estimating the likelihood of occurrence and magnitude of scenarios.
\item Explain steps in the process of defining and developing market-driven scenarios, and describe the application of the MDS approach in the example provided.
\item Describe challenges and considerations that may arise when implementing the market-driven scenario approach.
\end{enumerate}
\end{imkeybox}

% ------------------------------------------------------------------
\lo{IM-15 a}{Describe and evaluate the market-driven scenario (MDS) approach used to develop scenarios for stress testing investment portfolios.}

\begin{itemize}
\item \term{Historical replay} reruns a past episode such as 2008 or 1998. It is
defensible and easy to explain, but \term{the next crisis rarely takes the same
shape}, and the exposures in the portfolio today may have no counterpart in the
historical episode.
\item \term{Macro-narrative stress testing} starts from a path for GDP,
unemployment and inflation. It is the \term{regulatory norm}, but it requires a
long and uncertain mapping from macro variables down to the market prices that
actually move the portfolio.
\end{itemize}

\begin{imdefbox}[title={What MDS does instead}]
Shock a \emph{small set of core market variables directly}, then propagate those
shocks to every other asset using the estimated joint distribution of returns ---
so the scenario is expressed in the same units as the portfolio from the start,
with no macro-to-market mapping to defend.
\end{imdefbox}

% ------------------------------------------------------------------
\lo{IM-15 b}{Describe the scenario z-score, the volatility z-score, the correlation z-score, and the Mahalanobis distance, and explain how these are used in estimating the likelihood of occurrence and magnitude of scenarios.}

\renewcommand{\arraystretch}{1.25}
\noindent\begin{tabularx}{\textwidth}{@{}>{\raggedright\arraybackslash}p{2.3cm}p{2.6cm}X@{}}
\toprule
\textbf{\sffamily\small Measure} & \textbf{\sffamily\small Expression} & \textbf{\sffamily\small What it tells you} \\
\midrule
Scenario z-score & $z = (x - \mu)/\sigma$ & How extreme \emph{one} move is against its own history. A $z$ of 3 is a three-standard-deviation move. Larger absolute values mean more severe and less likely. \term{Univariate only.} \\
Volatility z-score & $z_\sigma = (\sigma_{\text{scen}} - \mu_\sigma)/\sigma_\sigma$ & Whether the scenario implies an unusually turbulent \emph{regime}. A scenario can use modest return shocks while still assuming a volatility environment that is itself historically extreme. \\
Correlation z-score & $z_\rho = (\rho_{\text{scen}} - \mu_\rho)/\sigma_\rho$ & Whether the assumed \emph{co-movement} is out of line with history --- for instance equities and bonds falling together when they normally offset. \\
Mahalanobis distance & $d^2 = (x-\mu)'\Sigma^{-1}(x-\mu)$ & The \term{multivariate} analogue: how far the whole shock vector sits from the centre of the joint distribution, \emph{accounting for the correlations between the shocks}. \\
\bottomrule
\end{tabularx}
\renewcommand{\arraystretch}{1}
\figcap{Three univariate measures and one multivariate one. The first three each
standardise a single feature of the scenario; only the Mahalanobis distance asks
whether the shocks are plausible \emph{together}.}

\figwrap{%
\begin{tikzpicture}
\begin{axis}[frmaxis,
  width=9.6cm, height=5.6cm,
  xmin=-4.2, xmax=4.2, ymin=-4.2, ymax=4.2,
  xlabel={shock to asset 1 (standard deviations)},
  ylabel={shock to asset 2},
  xtick={-3,0,3}, ytick={-3,0,3},
]
\draw[FRMRule, thick, rotate around={35:(axis cs:0,0)}]
  (axis cs:0,0) ellipse [x radius=3.3cm, y radius=1.0cm];
\draw[FRMRule, thin, rotate around={35:(axis cs:0,0)}]
  (axis cs:0,0) ellipse [x radius=2.1cm, y radius=0.62cm];
\addplot[only marks, mark=*, mark size=2.4pt, FRMGreen] coordinates {(2.4,2.4)};
\addplot[only marks, mark=*, mark size=2.4pt, FRMRed] coordinates {(2.4,-2.4)};
\draw[black!35, thin] (axis cs:-4.2,0) -- (axis cs:4.2,0);
\draw[black!35, thin] (axis cs:0,-4.2) -- (axis cs:0,4.2);
\node[font=\scriptsize\sffamily, fill=white, inner sep=1.5pt, anchor=south west,
      text=FRMGreen, anchor=south east] at (axis cs:2.25,2.55) {A: plausible};
\node[font=\scriptsize\sffamily, fill=white, inner sep=1.5pt, anchor=north west,
      text=FRMRed, anchor=north east] at (axis cs:2.25,-2.55) {B: implausible};
\node[font=\scriptsize\sffamily, fill=white, inner sep=1.5pt, anchor=west,
      text=black!55] at (axis cs:-4.0,3.6) {joint distribution};
\end{axis}
\end{tikzpicture}}%
{Why the univariate measures are not enough. \term{Points A and B have identical
z-scores on both axes} --- each is a $2.4$-sigma move in each asset. A sits along
the grain of the joint distribution and is plausible; B cuts against it and is far
less likely. \term{Only the multivariate measure separates them}, because only it
knows that the two assets normally move together.}

\begin{imkeybox}[title={How the two readings of a scenario combine}]
\term{Severity is read from the z-scores; likelihood is read from the Mahalanobis
distance.} A scenario is useful when it is \emph{severe enough to matter and close
enough to be defensible}. If $d$ is very large the scenario may be dismissed as
fantasy by the board; if $d$ is small it will not stress anything.
\end{imkeybox}

% ------------------------------------------------------------------
\lo{IM-15 c}{Explain steps in the process of defining and developing market-driven scenarios, and describe the application of the MDS approach in the example provided.}

\begin{enumerate}
\item \term{Define the narrative.} State the triggering event and its economic
logic --- a sovereign credit event, an energy shock, an abrupt policy surprise, a
geopolitical escalation. The narrative exists to make the scenario
\term{communicable} and to \term{discipline the choice of core assets}.
\item \term{Shock the core assets.} Identify the small set of variables the event
acts on directly and set the size and sign of each shock. Calibrate from
historical analogues, from option-implied moves, or from expert judgement.
\term{Keep the core small; every additional anchor is another judgement to
defend.}
\item \term{Test plausibility.} Compute the scenario, volatility and correlation
z-scores and the Mahalanobis distance for the core vector. If the combination is
implausible, revise the shocks or the narrative before going further. \term{This
is the step that distinguishes MDS from ad hoc scenario invention.}
\item \term{Propagate.} Map the core shocks onto every other exposure using the
estimated joint distribution, and revalue the portfolio.
\end{enumerate}

% ------------------------------------------------------------------
\lo{IM-15 d}{Describe challenges and considerations that may arise when implementing the market-driven scenario approach.}

\begin{itemize}
\item \term{Correlation instability and breakdown.} Diversification fails
precisely when the scenario bites. Conditional correlations \term{rise toward one
in a crisis}, so a matrix fitted on ordinary conditions \term{systematically
understates the loss}.
\item \term{Linearity of propagation.} The conditional expectation is a
\term{linear projection}. Options, convertibles, credit and anything with an
asymmetric payoff require \term{full revaluation}, and their own second-order
behaviour is not captured by the projection at all.
\item \term{Conditional mean is not a worst case.} Propagation gives the
\emph{expected} move in peripheral assets given the core shock. The residual
uncertainty around that expectation is discarded, so \term{a single point scenario
understates the dispersion of outcomes}.
\item \term{Plausibility against severity.} The two pull in opposite directions. A
scenario extreme enough to be interesting may be too implausible to be accepted.
\end{itemize}

\begin{imtrapbox}[title={Consolidated trap summary --- IM-15}]
\begin{itemize}
\item \term{Sibling, IM-15 a.} Historical replay reruns a past episode;
macro-narrative starts from GDP and inflation and is the \emph{regulatory norm};
MDS shocks market variables directly. Only MDS avoids the macro-to-market mapping.
\item \term{Scope, IM-15 b.} The three z-scores are \emph{univariate}; only the
Mahalanobis distance is multivariate. Two scenarios with identical z-scores can
have very different plausibility.
\item \term{Role, IM-15 b.} \term{z-scores give severity; Mahalanobis distance
gives likelihood.} Swapping the two is the standard corruption.
\item \term{Polarity, IM-15 d.} Correlations rise toward one in a crisis, so a
normal-period matrix \emph{understates} the loss. It does not overstate it.
\item \term{Definition, IM-15 d.} Propagation returns a \emph{conditional mean},
not a worst case. The dispersion around it is thrown away.
\end{itemize}
\end{imtrapbox}

% ==================================================================
